% !TEX root = ../main.tex

\chapter{C Programming}

\section{How a Computer Runs a Program}

Before writing C code, students should understand what a program does inside a computer. A program receives data when necessary, gives instructions for processing it, and produces a result.

\begin{definitionbox}[Program]
    A \term{program} is an ordered set of instructions written to make a computer perform a specific task or solve a problem.
\end{definitionbox}

\begin{definitionbox}[Programming Language]
    A \term{programming language} is a formal language used to write the instructions of a program in a form that can be translated and executed by a computer.
\end{definitionbox}

\subsection{Input, Processing and Output}

Every useful program follows an information-processing path: it receives input, processes it according to the program, and produces output. Input can come from a keyboard, file, sensor or another device; output can appear on a screen, printer, file or another device.

\begin{formulabox}[Information-Processing Path]
    \centering
    \term{Input} \(\longrightarrow\) \term{Processing by CPU} \(\longrightarrow\) \term{Output}
\end{formulabox}

For example, when a student enters two numbers through a keyboard, a program can instruct the computer to add them. The entered numbers are the input, addition is the processing, and the displayed sum is the output.

\subsection{Role of the CPU While a Program Runs}

The \term{Central Processing Unit (CPU)}, or processor, executes program instructions and processes data. Its main parts work together as follows:

\begin{itemize}
    \item \term{Control Unit (CU)}: selects instructions in order and directs the required operation.
    \item \term{Arithmetic and Logic Unit (ALU)}: performs calculations and logical tests such as comparisons.
    \item \term{Registers / working memory}: temporarily hold instructions, data and intermediate results.
\end{itemize}

The execution path can therefore be understood as follows:

\begin{enumerate}
    \item Input data and the required instructions are made available to the CPU.
    \item The control unit selects the next instruction and directs the required operation.
    \item The ALU performs a calculation or logical test when the instruction requires one.
    \item Intermediate and final results are held temporarily while processing continues.
    \item The final result is sent to an output device or stored for later use.
\end{enumerate}

